\section{Dimension-dependent output fluctuations}
The local tail also identifies why convergence at dimension four can be
slow. Gaussian and stable condensate fluctuations are established
phenomena in heavy-tailed conditioned sums; see, for example,
\href{https://arxiv.org/abs/0912.1793}{Armend\'ariz--Grosskinsky--Loulakis},
Section 2.3, including its discussion of infinite-variance bulk laws.
The result here computes the dimension threshold, spectral measure, and
output projection for the Gaussian ReLU statistic. It does not claim
a new general stable-limit principle. Let $f_s=\sigma_s\phi_{v_s}$, so
$\langle\chi,f_s\rangle_{L^2(X)}=\rho$, and define the projection
\[
 \mathcal P_s y=y-\rho^{-1}f_s\langle\chi,y\rangle_{L^2(X)}.
\]
This need not be an orthogonal projection. It removes the one scalar
component constrained by the conditioning event.

\begin{theorem}[Stable and Gaussian fluctuation regimes]
Keep fixed $d\ge4$ and even $2\le r\le d$, and $s>m_c$.
Let $S_N^{\rm ray}$ be the nearest critical orientation of the largest
neuron, assigning any value off the event where it is well defined.
Put $\alpha=(d-1)/2$ and
\[
 b_N=\begin{cases}
 N^{2/3},&d=4,\\
 \sqrt{N\log N},&d=5,\\
 \sqrt N,&d\ge6.
 \end{cases}
\]
Under scalar conditioning, the following joint convergence holds:
\[
 \left(S_N^{\rm ray},\frac N{b_N}
    \left[F_N-m_c\chi-\frac\delta\rho f_{S_N^{\rm ray}}\right]\right)
       \Longrightarrow (S,\mathcal P_S Y),                \tag{21}
\]
where $S$ is uniform among the critical orientations and independent
of $Y$. The distribution of $Y$ is as follows.
\begin{enumerate}
\item For $d=4$, it is the centered $3/2$-stable vector with
characteristic exponent
\[
 \log\E e^{i\xi\cdot Y}
 =\frac{\alpha C}{2^r}\sum_s\int_0^\infty
  \left(e^{it\xi\cdot f_s/\rho}-1-it\xi\cdot f_s/\rho\right)
                           t^{-\alpha-1}\,dt.            \tag{22}
\]
Here $\xi\cdot y$ denotes the ordinary Euclidean pairing on the
finite output-vector space.
\item For $d=5$, it is centered Gaussian with covariance
\[
                 \Gamma=\frac{C}{2^r\rho^2}
                                   \sum_s f_sf_s^\top.   \tag{23}
\]
\item For $d\ge6$, it is centered Gaussian with covariance
$\Sigma=\operatorname{Cov}_{Q_*}(\Phi)$, which is finite.
\end{enumerate}
The stable or Gaussian vector can be degenerate on some output
directions. No finite-width error bound is asserted.
\end{theorem}
\begin{proof}
We first identify the independent critical-bulk limit. The density and
mark lemma implies multivariate regular variation of $\Phi$ under $Q_*$:
on sets bounded away from zero its tail measure is
\[
 \nu(B)=\frac{\alpha C}{2^r}\sum_s\int_0^\infty
                    \mathbf1_{\{tf_s/\rho\in B\}}
                                      t^{-\alpha-1}\,dt. \tag{24}
\]
To justify that no additional output tails are omitted, write
$D_0=\|\theta\|^2/2-Z/\rho$. Outside any fixed conical neighborhood
of the critical rays, $D_0\ge c\|\theta\|^2$ by compactness of the
unit sphere. The same gap holds on
$Z\le\rho\|\theta\|^2/4$. Since $\|\Phi\|\le C_d\|\theta\|^2$,
large output on either of these sets has an exponentially small tail.
On the remaining ray tubes, the local coordinate calculation gives
$\Phi/Z\to f_s/\rho$ with the stated uniform ray weights. It also
gives the bound $Q_*(\|\Phi\|>t)\le C't^{-\alpha}$ for large $t$.
These facts prove (24) and the required truncated-moment bounds.

For $d=4$, apply (24) to
$e^{i\xi\cdot y}-1-i\xi\cdot y$ with scale $b_N=N^{1/\alpha}$.
The contribution from $\|y\|\le\eta$ is bounded by $O(\eta^{2-\alpha})$
using the truncated second moment. The contribution from
$\|y\|>R$ is bounded by $O(R^{1-\alpha})$ using the tail first moment.
On the intervening annulus, tail-measure convergence applies directly.
Thus the characteristic exponent of
$b_N^{-1}\sum_{j<N}(\Phi_j-m_c\chi)$ tends to (22).
The quadratic error in expanding the logarithm is negligible because
$N/b_N^2\to0$ and the first moment is finite. This proves the stable
bulk limit without assuming a finite variance.

For $d=5$, (24) with $\alpha=2$ instead gives
\[
 \E[\Phi\Phi^\top\mathbf1_{\{\|\Phi\|\le R\}}]
       \sim \frac{2C\log R}{2^r\rho^2}\sum_s f_sf_s^\top
\]
in the sense of entrywise leading logarithmic coefficients; bounded
centering changes no coefficient. With $b_N=\sqrt{N\log N}$,
$\log b_N/\log N\to1/2$, giving (23).
The chance of a summand above $b_N$ is $O(1/\log N)$ after a union
bound. The centered mean removed by this truncation, divided by $b_N$,
is also $O(1/\log N)$. For each fixed $\eta>0$, the variance of
summands between $\eta b_N$ and $b_N$, divided by $b_N^2/N$, tends
to zero. The elementary triangular-array Lindeberg theorem therefore
gives the claimed Gaussian bulk limit. For $d\ge6$, the ordinary
finite-variance multivariate central limit theorem applies directly.

Finally the joint product law proved from (15a) makes these bulk limits
valid under the conditioned law, independently of the exceptional
orientation. Its tight mark law gives
\[
 \Phi(\theta_{J_N})=\rho^{-1}Z_{J_N}f_{S_N^{\rm ray}}+O_P(1).
\]
Indeed $\Phi=c f_s+O_P(1)$ and $Z=\rho c+O_P(1)$ in the tight
$(v,t)$ coordinates. Since
$Z_{J_N}=Ns+u-\sum_{i\ne J_N}Z_i$ with tight $u$, substitution
into the total output gives exactly $\mathcal P_{S_N^{\rm ray}}$
applied to the centered bulk sum, up to $O_P(1)$.
Dividing by $b_N\to\infty$ proves (21).
\end{proof}

In particular dimension-four output fluctuations are of order
$N^{-1/3}$, rather than $N^{-1/2}$. The logarithmic boundary at
dimension five and the finite-variance regime above it concern the
critical ordinary neurons. They do not follow from a finite-variance
approximation imposed on the exceptional neuron.
